% =============================================================================================
%  EpistemeGraph -- manuscript
%  Every quantitative result is a macro defined in generated/results.tex, which is written by
%  scripts/export_paper_assets.py from the actual run artefacts. Un-run results render as "--".
%  Compile:  latexmk -pdf main.tex      (or pdflatex -> bibtex -> pdflatex x2)
%  Overleaf: upload paper/EpistemeGraph_overleaf.zip (New Project -> Upload Project) and press Recompile.
%            The zip is produced by scripts/export_paper_assets.py or the website's "Export Overleaf bundle".
% =============================================================================================
\documentclass[10pt,twocolumn]{article}

\usepackage[margin=0.75in,columnsep=0.28in]{geometry}
\usepackage[T1]{fontenc}
\usepackage[utf8]{inputenc}
\usepackage{lmodern}
\usepackage{microtype}
\usepackage{amsmath,amssymb,amsthm,mathtools}
\usepackage{bm}
\usepackage{booktabs}
\usepackage{graphicx}
\usepackage{xcolor}
\usepackage{algorithm}
\usepackage{algpseudocode}
\usepackage[numbers,sort&compress]{natbib}
\usepackage[colorlinks=true,linkcolor=blue!50!black,citecolor=blue!50!black,urlcolor=blue!50!black]{hyperref}
\usepackage[capitalise,noabbrev]{cleveref}

\newtheorem{definition}{Definition}
\newtheorem{proposition}{Proposition}
\newtheorem{lemma}{Lemma}
\theoremstyle{remark}
\newtheorem{remark}{Remark}

\DeclareMathOperator{\logit}{logit}
\DeclareMathOperator{\spec}{spec}
\DeclareMathOperator*{\argmin}{arg\,min}
\newcommand{\R}{\mathbb{R}}
\newcommand{\E}{\mathbb{E}}
\newcommand{\SKGq}{S_{\mathrm{KG}}(q)}
\newcommand{\todo}[1]{\textcolor{red!70!black}{[#1]}}

\IfFileExists{generated/results.tex}{\input{generated/results.tex}}{%
  \newcommand{\SKG}{--}\newcommand{\PhiProv}{--}\newcommand{\PhiCov}{--}\newcommand{\PhiTopo}{--}\newcommand{\PhiCons}{--}
  \newcommand{\NDocs}{--}\newcommand{\NChunks}{--}\newcommand{\NRawTriples}{--}\newcommand{\NTriples}{--}
  \newcommand{\NEntities}{--}\newcommand{\NShared}{--}\newcommand{\AcceptRate}{--}\newcommand{\NItems}{--}
  \newcommand{\NLabelled}{--}\newcommand{\NTrain}{--}\newcommand{\NCal}{--}\newcommand{\NTest}{--}
  \newcommand{\SupportedRate}{--}\newcommand{\Champion}{--}\newcommand{\TestAUC}{--}\newcommand{\TestLL}{--}
  \newcommand{\TestBrier}{--}\newcommand{\TestECE}{--}\newcommand{\TestAcc}{--}\newcommand{\LLMAUC}{--}
  \newcommand{\LLMECE}{--}\newcommand{\Alpha}{--}\newcommand{\Coverage}{--}\newcommand{\CoverageSup}{--}
  \newcommand{\CoverageHal}{--}\newcommand{\AbstainRate}{--}\newcommand{\SelAcc}{--}\newcommand{\TauVer}{--}
  \newcommand{\TauHal}{--}\newcommand{\PoolBzero}{--}\newcommand{\PoolBq}{--}\newcommand{\PoolBkg}{--}
  \newcommand{\JudgeKappa}{--}\newcommand{\JudgeN}{--}
  \newcommand{\PrimaryJudge}{--}\newcommand{\CrossJudge}{--}\newcommand{\CrossJudgeN}{--}\newcommand{\CrossJudgeKappa}{--}\newcommand{\CrossJudgeAgree}{--}
  \newcommand{\KGPrec}{--}\newcommand{\KGPrecLo}{--}\newcommand{\KGPrecHi}{--}\newcommand{\KGPrecLenient}{--}
  \newcommand{\KGAuditN}{--}\newcommand{\KGAuditor}{--}\newcommand{\ExtentHal}{--}\newcommand{\ExtentSup}{--}\newcommand{\ExtentAUC}{--}
  \newcommand{\TestAUCLo}{--}\newcommand{\TestAUCHi}{--}\newcommand{\LLMAUCLo}{--}\newcommand{\LLMAUCHi}{--}
  \newcommand{\DeltaLLM}{--}\newcommand{\PvalLLM}{--}\newcommand{\DeltaText}{--}\newcommand{\PvalText}{--}
  \newcommand{\DeltaTopo}{--}\newcommand{\PvalTopo}{--}\newcommand{\KGAccLearned}{--}\newcommand{\KGAccLo}{--}\newcommand{\KGAccHi}{--}
  \newcommand{\SensDefault}{--}\newcommand{\SensMaxDev}{--}\newcommand{\SensRange}{--}
  \newcommand{\TableLeaderboard}{}\newcommand{\TableAblation}{}\newcommand{\TableFeatures}{}\newcommand{\TablePercat}{}\newcommand{\TableKgprec}{}
  \newcommand{\TableSens}{}\newcommand{\TableTripleweights}{}}

\title{\textbf{EpistemeGraph: Knowledge-Graph-Grounded Hallucination Auditing\\with Sheaf, Curvature and Spectral Evidence and Conformal Guarantees}}
\author{Anurag Raj\\
  \small Indian Institute of Technology \todo{department, campus}\\
  \small \texttt{\todo{email}}\quad \url{https://github.com/Virtuoso-8051/EpistemeGraph}}
\date{}

\begin{document}
\maketitle

% =============================================================================================
\begin{abstract}
Large language models (LLMs) answer questions about specialised scientific literature fluently, and with high
verbalised confidence, even when they invent mechanisms, misattribute results across papers or accept a false
premise. We present \emph{EpistemeGraph}, an auditing system that decides whether an LLM answer is supported by a
document collection, how likely it is to be hallucinated, how much of it is unsupported, and which claims are at
fault. EpistemeGraph (i)~builds a provenance-verified knowledge graph (KG) from PDFs and measures it twice: with a
decomposable intrinsic fidelity index $S_{\mathrm{KG}}$, and with a \emph{learned}, model-assisted estimate of KG
accuracy with a confidence interval; (ii)~converts each answer into a polarity-annotated claim graph through a call
decoupled from answer generation; (iii)~describes the answer by a 13-dimensional vector that combines three
topological measurements of the claim graph against the KG (the Dirichlet energy of a translational knowledge sheaf,
the Wasserstein shift of the Forman--Ricci curvature distribution when the claims are grafted onto the KG, and a
spectral divergence of normalised Laplacians) with graph provenance, natural-language-inference consistency,
out-of-corpus novelty and model-side uncertainty, and by a 15-dimensional vector per claim; and (iv)~learns the
decision end to end, with no hand-set weight or threshold: a ten-model nested-CV tournament, including an
attention-based multiple-instance network that localises unsupported claims from answer-level labels only, selects a
calibrated scorer; a learned log-linear opinion pool adds the local KG reliability; and Mondrian split-conformal
prediction returns \textsc{verified}, \textsc{abstain} or \textsc{hallucinated} with class-conditional coverage. On a
\NItems-question benchmark over the LINE, node2vec and DeepWalk papers (factual, cross-document, unanswerable and
false-premise questions, labelled at the answer level), the audit reaches AUROC \TestAUC{}
(95\% CI [\TestAUCLo, \TestAUCHi]) on held-out answers against \LLMAUC{} for the LLM's own confidence (DeLong
$p$ = \PvalLLM), with empirical conformal coverage of \Coverage\% at $\alpha=\Alpha$. The KG fidelity index is stable
under perturbation of every Phase-1 design constant (maximum change \SensMaxDev). Code, data, prompts and every
intermediate artefact are public.
\end{abstract}

% =============================================================================================
\section{Introduction}
\label{sec:intro}

Hallucination, i.e.\ content that is fluent but unsupported by the relevant source, is the central obstacle to
deploying LLMs in technical domains~\cite{ji2023survey,huang2025survey}. Two facts make it hard to manage in practice.
First, LLM self-assessment is an unreliable safety signal: verbalised confidence is typically high and weakly
discriminative on specialised material~\cite{xiong2024uncertainty,kadavath2022know}. Second, most reference-based
verifiers compare text with text, which cannot distinguish a claim that \emph{uses} the corpus vocabulary from one that
is \emph{supported} by the corpus. For example, ``node2vec optimises a genetic fitness function'' shares many tokens
with the node2vec paper and is still false.

Knowledge graphs offer a structured alternative~\cite{kim2023factkg,sansford2024grapheval}. If the source documents are
distilled into a graph of typed relations with evidence, an answer can be distilled into a claim graph and the two
graphs compared. This raises four questions that we address directly.
\textbf{(Q1)} The reference graph is itself LLM-extracted. How accurate is it, measured rather than assumed, and how
should its reliability enter the decision? \textbf{(Q2)} Which quantities expose an unsupported claim beyond counting
missing edges? \textbf{(Q3)} How should many such signals be turned into a decision without hand-tuned weights and
thresholds, and with a statistical guarantee? \textbf{(Q4)} Can the system say \emph{which} claim is wrong when it is
only ever told whether a whole answer is wrong?

\paragraph{Contributions.}
\begin{enumerate}\itemsep0pt
\item \textbf{Measured KG quality.} A provenance-verified, corpus-namespaced KG construction with an intrinsic
  fidelity index $S_{\mathrm{KG}}$ whose aggregation rule is derived, not tuned (\cref{prop:gm}); an identifiability
  argument for why its weights cannot be learned from one corpus; and a \emph{learned} model-assisted estimator of KG
  accuracy that is design-unbiased and tighter than the plain sample proportion (\cref{sec:phase1}).
\item \textbf{Decoupled interception}, so the audited model cannot omit its own fabrications from the claim graph
  (\cref{sec:phase2}).
\item \textbf{Geometric evidence.} Three topological measurements of a claim graph against a reference graph,
  with stated properties: knowledge-sheaf energy, Forman--Ricci curvature shift under grafting (zero exactly for
  answers that add nothing new), and a bounded spectral divergence against the reference's $k$-hop shadow. They are
  combined with graph provenance, NLI consistency and model-side uncertainty, all computed on the \emph{answer}
  (\cref{sec:phase3}).
\item \textbf{A fully learned decision layer} with claim-level localisation: a nested-CV tournament under a proper
  scoring rule whose contenders include a gated-attention multiple-instance network trained with a noisy-OR
  constraint; a log-linear pool that strictly nests fixed weightings; and Mondrian conformal prediction for a
  three-way decision with class-conditional coverage (\cref{sec:phase4}).
\item \textbf{Evaluation protocol.} A \NItems-question benchmark with evidence-validated items and answer-level
  labels, bootstrap confidence intervals, DeLong tests against the LLM's own confidence and against a text-only NLI
  checker, family ablations, and a robustness analysis of every Phase-1 design constant
  (\cref{sec:setup,sec:results}).
\end{enumerate}

\paragraph{What is new, and what is not.}
The individual instruments are established: cellular and knowledge sheaves~\cite{hansen2019sheaves,gebhart2023knowledge},
Forman curvature~\cite{forman2003bochner,samal2018comparative}, Laplacian spectra~\cite{chung1997spectral}, NLI-based
consistency~\cite{laban2022summac}, attention-based multiple-instance learning~\cite{ilse2018attention}, log-linear
pooling~\cite{genest1986combining}, model-assisted survey estimation~\cite{sarndal1992model} and conformal
prediction~\cite{vovk2005algorithmic,vovk2012conditional}. Our contribution is (a)~their use as \emph{measurement
instruments} for the agreement between an answer's claim graph and a fallible reference graph, with new definitions
for the two graft-based quantities ($\Delta\kappa$, $D_{\mathrm{spec}}$) and their properties
(\cref{prop:curv,lem:spec}); (b)~the identifiability analysis that makes the KG reliability a learnable,
per-question quantity (\cref{prop:pool}); (c)~treating an answer as a bag of claims, so that answer-level labels
supervise claim-level hallucination localisation; and (d)~an end-to-end protocol in which every weight and threshold
is estimated from data and every remaining design constant is shown not to drive the conclusions.

% =============================================================================================
\section{Related Work}
\label{sec:related}

\paragraph{Hallucination detection.}
Reference-free detectors use the model itself: sampling consistency~\cite{manakul2023selfcheckgpt}, semantic
entropy over meaning clusters~\cite{farquhar2024semantic}, or elicited confidence~\cite{kadavath2022know,xiong2024uncertainty}.
Reference-based methods decompose text into atomic facts and verify them against a knowledge
source~\cite{min2023factscore}. Benchmarks such as HaluEval~\cite{li2023halueval} and RAGTruth~\cite{niu2024ragtruth}
established response-level and span-level annotation. Our labels follow the response-level convention, and our
false-premise and unanswerable categories target the failure modes most relevant to literature QA.

\paragraph{Graphs for grounding and verification.}
Retrieval augmentation~\cite{lewis2020rag} and graph-structured retrieval~\cite{edge2024graphrag} reduce
hallucination at generation time. FactKG~\cite{kim2023factkg} verifies claims by reasoning over a KG, and
GraphEval~\cite{sansford2024grapheval} represents an output as a KG and checks each triple with an NLI model.
EpistemeGraph differs in three ways: it treats the reference graph as fallible and measures its quality; it uses
graph-\emph{geometric} rather than purely per-triple evidence; and it attaches a distribution-free guarantee to the
final decision.

\paragraph{Sheaves, curvature and spectra.}
Cellular sheaves generalise graph Laplacians by attaching vector spaces and linear restriction maps to vertices and
edges~\cite{hansen2019sheaves}. They underlie sheaf neural networks~\cite{bodnar2022sheaf} and knowledge-graph
embedding, where TransE~\cite{bordes2013transe} is a special case of a knowledge sheaf~\cite{gebhart2023knowledge}.
Forman's combinatorial Ricci curvature~\cite{forman2003bochner} and its network
adaptations~\cite{sreejith2016forman,samal2018comparative} identify bridge-like edges, which are linked to
bottlenecks in message passing~\cite{topping2022oversquashing}. Laplacian spectra summarise global
shape~\cite{chung1997spectral}. We use all three as measurement instruments for claim graphs, not as learning
architectures.

\paragraph{Calibration and conformal prediction.}
Proper scoring rules~\cite{gneiting2007scoring} and post-hoc calibration~\cite{platt1999probabilistic,guo2017calibration}
make scores interpretable as probabilities. Conformal prediction~\cite{vovk2005algorithmic,lei2018distribution,angelopoulos2023gentle}
gives finite-sample coverage under exchangeability, and its Mondrian variant gives class-conditional
validity~\cite{vovk2012conditional}. Pooling probabilities from several sources has a long history in
statistics~\cite{genest1986combining}. We use logarithmic pooling with learned weights.

\paragraph{Corpus.}
Our experimental corpus consists of three representation-learning papers: DeepWalk~\cite{perozzi2014deepwalk},
LINE~\cite{tang2015line} and node2vec~\cite{grover2016node2vec}. Their overlapping vocabulary (random walks, proximity,
Skip-gram, negative sampling) with different mechanisms creates realistic opportunities for cross-paper
confabulation.

% =============================================================================================
\section{Problem Formulation}
\label{sec:problem}

Let $\mathcal D$ be a document collection and $q$ a user query. An LLM $M$ returns an answer $a=M(q)$ and a verbalised
confidence $c_{\mathrm{LLM}}\in[0,1]$. Let $\mathcal C^{+}(a)$ be the set of factual claims that $a$ asserts. We call
$a$ \emph{supported} ($y=1$) if every claim in $\mathcal C^{+}(a)$ is stated in or entailed by $\mathcal D$ (or is
uncontroversial background), and \emph{hallucinated} ($y=0$) otherwise. Correctly rejecting a false premise, or
stating that $\mathcal D$ does not contain an answer, is supported behaviour.

\emph{Goal:} a map $(q,a)\mapsto\Gamma\subseteq\{0,1\}$ such that
$\Pr\big(y\in\Gamma\mid y=c\big)\ge 1-\alpha$ for each class $c$, with $|\Gamma|=1$ as often as possible. The
three-way decision follows from $\Gamma$: $\{1\}\mapsto$\textsc{verified}, $\{0\}\mapsto$\textsc{hallucinated},
otherwise \textsc{abstain}.

% =============================================================================================
\section{Method}
\label{sec:method}

\subsection{Phase 1: Incremental KG Construction and Fidelity}
\label{sec:phase1}

\paragraph{Extraction with verified provenance.}
Each PDF receives a content hash. Its text is normalised (NFKC, de-hyphenation, bibliography removal) and packed into
sentence-aligned chunks of about 2{,}200 characters with one sentence of overlap. An LLM extracts entities and then
triples whose endpoints must come from the entity list. Each triple carries a verbatim evidence sentence. A triple
is \emph{accepted} only if (a)~its quote matches the chunk verbatim (fuzzy ratio $\ge0.82$) or anchors to a chunk
sentence by embedding similarity, and (b)~both entities occur in that sentence. Every extracted triple $c$ receives a
grounding score $g_c\in[0,1]$, the product of an evidence-match score and a blend of entity support and semantic
support. Rejected triples are logged.

\paragraph{Incremental entity resolution and storage.}
Mentions are resolved to canonical entities by exact match, acronym$\leftrightarrow$expansion match, and type-aware
cosine similarity $\ge0.85$ on sentence embeddings~\cite{reimers2019sbert}. When new documents arrive, the resolver
is \emph{pre-seeded} with all existing canonical entities, so a new surface form merges into an existing node and
existing names are never renamed. The delta is written to Neo4j with \texttt{MERGE} keyed on a normalised name, with
list properties unioned and evidence replaced only by better-grounded evidence (\cref{app:cypher}). Chunk identifiers
are content-addressed, which makes ingestion idempotent and keeps support counts exact.

\paragraph{Fidelity $S_{\mathrm{KG}}$.}
Let $\mathcal C$ be all extracted triples, $\mathcal C^{+}\subseteq\mathcal C$ the accepted ones, $G=(V,E)$ the
resulting graph and $\mathcal K$ the chunks. We define four pillars in $[0,1]$:
\begin{align}
\Phi_{\mathrm{prov}} &= \tfrac{1}{|\mathcal C|}\textstyle\sum_{c\in\mathcal C} g_c, \\
\Phi_{\mathrm{cov}}  &= \frac{\sum_{k} w_k\,\mathrm{cov}(k)}{\sum_k w_k},\quad w_k=\log(1+|k|),\\
\Phi_{\mathrm{topo}} &= 1-\frac{\kappa(G)-1}{|V|-1}, \\
\Phi_{\mathrm{cons}} &= 1-\frac{\#\{\text{functional violations}\}}{|\mathcal C^{+}|},
\end{align}
where $\mathrm{cov}(k)=\frac{1}{m_k}\sum_{i}\min\!\big(1,\max_{t\in T_k}\cos(\phi(s_i),\phi(t))/\tau\big)$ measures how
well the triples $T_k$ extracted from chunk $k$ cover its $m_k$ sentences, and $\kappa(G)$ counts connected
components. $\Phi_{\mathrm{prov}}$ is a precision proxy, $\Phi_{\mathrm{cov}}$ a recall proxy, $\Phi_{\mathrm{topo}}$
equals $1$ for a connected graph and $0$ for isolated triples, and $\Phi_{\mathrm{cons}}$ penalises a subject with two
objects under a functional relation. We aggregate with the unweighted geometric mean
\begin{equation}
S_{\mathrm{KG}}=\big(\Phi_{\mathrm{prov}}\Phi_{\mathrm{cov}}\Phi_{\mathrm{topo}}\Phi_{\mathrm{cons}}\big)^{1/4},
\label{eq:skg}
\end{equation}
the symmetric, scale-invariant mean that vanishes if any pillar does. All pillars are sums over chunks and triples,
so $S_{\mathrm{KG}}$ is recomputed exactly over the whole corpus after every incremental upload without
re-processing earlier documents.

\paragraph{Why the exponent $1/4$ is derived, not tuned.}
Consider the weighted power means $M_p^{\bm w}(\Phi)=\big(\sum_j w_j\Phi_j^{\,p}\big)^{1/p}$ ($M_0^{\bm w}=\prod_j\Phi_j^{w_j}$),
the family of homogeneous quasi-arithmetic means~\cite{hardy1934inequalities}. We require three properties of a
quality index: \emph{symmetry} (no pillar is privileged a priori, hence $w_j=1/4$), \emph{homogeneity}
(rescaling every pillar rescales the index), and \emph{veto} (a failed pillar cannot be compensated: $\Phi_j=0\Rightarrow
S=0$).
\begin{proposition}\label{prop:gm}
Among symmetric power means $M_p$ with the veto property, the geometric mean $M_0$ is the pointwise largest, i.e.\ the
least pessimistic aggregator that still vetoes.
\end{proposition}
\begin{proof}
For $p>0$, $M_p(\Phi)\ge(\tfrac14\Phi_j^{\,p})^{1/p}>0$ whenever some $\Phi_j>0$, so the veto fails. For $p\le0$ the veto
holds (with the convention $0^{p}=\infty$ for $p<0$). By the power-mean inequality $M_p$ is non-decreasing in
$p$~\cite{hardy1934inequalities}, so $\sup_{p\le0}M_p=M_0$.
\end{proof}
Learning the weights instead is not possible at this level: $S_{\mathrm{KG}}$ is one number per corpus, so with one
corpus there is one observation and the weights are unidentifiable (the same argument as \cref{prop:pool}(b)). The
learnable quantity lives one level down, at the level of individual triples, which we exploit next.

\paragraph{Measured, learned accuracy of the KG.}
$S_{\mathrm{KG}}$ is an index; it does not state how often a stored triple is \emph{correct}. We therefore audit a
uniform random sample of $n$ of the $N$ accepted triples: an auditor sees each triple with its full source passage
and labels it \textsc{correct}, \textsc{partially correct} or \textsc{incorrect}. With $k$ correct triples,
$\hat p=k/n$ has the 95\% Wilson interval~\cite{wilson1927probable} ($z=1.96$)
\begin{equation}
\mathrm{CI}_{95}(\hat p)=\frac{\hat p+\frac{z^2}{2n}\pm z\sqrt{\frac{\hat p(1-\hat p)}{n}+\frac{z^2}{4n^2}}}{1+\frac{z^2}{n}}.
\label{eq:wilson}
\end{equation}
The unaudited $N-n$ triples carry information too: every triple has evidence signals $\bm v_t$ (grounding score,
verbatim match, extraction confidence, support count, endpoint degrees, relation frequency, cross-paper status,
quote length). We learn $\hat m(\bm v)=\hat P(\text{correct}\mid\bm v)$ with a regularised logistic model on the audited
sample and use the model-assisted difference estimator of survey sampling~\cite{sarndal1992model}
\begin{equation}
\hat p_{\mathrm{MA}}=\frac1N\sum_{t=1}^{N}\hat m(\bm v_t)+\frac1n\sum_{t\in s}\big(y_t-\hat m^{(-t)}(\bm v_t)\big),
\qquad
\widehat{\mathrm{SE}}=\sqrt{\Big(1-\frac nN\Big)\frac{\widehat{\mathrm{Var}}(y_t-\hat m^{(-t)}(\bm v_t))}{n}},
\label{eq:ma}
\end{equation}
where $\hat m^{(-t)}$ are cross-fitted predictions. Under simple random sampling $\hat p_{\mathrm{MA}}$ is unbiased for
the KG accuracy \emph{whatever} the model, and its variance falls as the model explains more of the label variance:
the learned weights (\cref{tab:tripleweights}) buy precision, never bias. On our corpus the learned accuracy is
\KGAccLearned\% (95\% CI [\KGAccLo\%, \KGAccHi\%]).

\paragraph{Robustness of the remaining constants.}
Phase~1 still contains design constants that are not mathematical definitions (quote-match ratio, entity-support
level, coverage normaliser $\tau$, minimum extractor confidence, entity-merge cosine). Instead of tuning them, we
perturb each one at a time and re-run verification, entity resolution and $S_{\mathrm{KG}}$ on the cached extractions
(\cref{tab:sens}). $S_{\mathrm{KG}}$ changes by at most \SensMaxDev{} (default \SensDefault). The most influential
constant is $\tau$, which by construction rescales $\Phi_{\mathrm{cov}}$; the quote-match ratio and the extractor
confidence threshold have essentially no effect, because accepted quotes are near-verbatim and extractor confidences
are concentrated above both thresholds.

\paragraph{Corpus workspaces.}
Every uploaded set of documents defines a \emph{workspace}: a corpus with its own graph, written to Neo4j under a
dedicated namespace (a \texttt{kg\_id} property with a composite uniqueness constraint on
$(\texttt{kg\_id},\texttt{name\_key})$). Uploading a new set therefore builds a fresh graph from scratch without
modifying earlier ones, while an explicit \emph{extend} operation merges new documents into an existing corpus.
Because the decision layer (\cref{sec:phase4}) is calibrated per corpus, the workspace is also the unit of
calibration.

\subsection{Phase 2: Decoupled Interception}
\label{sec:phase2}

The audited model first answers $q$ as it would in deployment and states $c_{\mathrm{LLM}}$ (\emph{closed-book}
mode; a \emph{grounded} mode adds retrieved corpus sentences and audits a RAG system). A \emph{separate} extraction
call then converts $a$ into the claim graph $G'=(V',E')$. Each edge carries a relation type and a polarity in
$\{$\textsc{asserted}, \textsc{speculative}, \textsc{refuted}$\}$. Decoupling matters: if a single call answers
\emph{and} lists its claims, the model controls which of its statements are audited. The extraction prompt is
assembled dynamically. Its relation vocabulary is mined from the live KG, and its naming hints are the KG entities
nearest to $q\oplus a$ in embedding space. All calls are cached by a hash of prompt version, model and input.

\subsection{Phase 3: Retrieval, Alignment and the CASSE Vector}
\label{sec:phase3}

\paragraph{Retrieval and alignment.}
Seed vertices are KG entities matched by name, alias, acronym or a distinctive token of the claim entities. The
ground-truth (GT) subgraph $G=(V,E)$ contains the seeds, their neighbours, and every shortest path of length $\le3$
between two seeds. Each response entity is aligned to at most one GT vertex (exact form, acronym, then cosine $\ge0.85$,
the Phase~1 merge threshold). Let $E'_{+}$ be the asserted and speculative claim edges and
$x_v=\phi(\mathrm{name}(v))\in\R^{d}$ unit embeddings. All features are computed on the response, and none uses the
query text.

\paragraph{(1) Sheaf Laplacian energy.}
Consider the translational knowledge sheaf $\mathcal F$ on a graph whose vertex and edge stalks are $\R^d$, with
identity restriction maps and a relation-specific affine shift $t_r$. Its coboundary on $e=(u\xrightarrow{r}v)$ is
$(\delta x)_e=x_v-x_u-t_r$, and $\sum_e\|(\delta x)_e\|^2$ is the Dirichlet energy of the affine sheaf
Laplacian~\cite{hansen2019sheaves,gebhart2023knowledge}. TransE~\cite{bordes2013transe} is precisely the model that
makes this energy small on true triples. We estimate $t_r$ on the GT subgraph by maximum likelihood under an
isotropic Gaussian residual, shrunk towards the global mean shift $\bar t$ with one pseudo-observation:
\begin{equation}
t_r=\frac{\sum_{(u,v)\in E_r}(x_v-x_u)+\bar t}{|E_r|+1}.
\end{equation}
The feature is the response's mean edge energy relative to the GT baseline:
\begin{equation}
\mathcal E_{\mathcal F}=\frac{\frac1{|E'_+|}\sum_{e\in E'_+}\|x_v-x_u-t_r\|^2}{\frac1{|E|}\sum_{e\in E}\|x_v-x_u-t_r\|^2}.
\end{equation}
$\mathcal E_{\mathcal F}=1$ means the claims are as relationally consistent as the ground truth. Large values indicate
relations whose endpoints do not stand in the semantic offset that the relation usually induces.

\paragraph{(2) Forman--Ricci curvature shift.}
For an unweighted graph, the augmented Forman curvature of an edge is~\cite{samal2018comparative}
\begin{equation}
F(u,v)=4-d_u-d_v+3\,|\triangle(u,v)|,
\end{equation}
where $|\triangle(u,v)|$ counts the triangles containing the edge. Bridges between otherwise unrelated regions are
strongly negatively curved, while edges with shared context gain curvature through triangles. We \emph{graft} the
response onto the ground truth and measure how far the curvature distribution moves:
\begin{equation}
\Delta\kappa=\log\Big(1+W_1\big(\mu_{G},\,\mu_{G\cup E'_+}\big)\Big),
\end{equation}
where $\mu_H$ is the empirical distribution of $\{F_H(e)\}_{e\in E(H)}$ and $W_1$ the 1-Wasserstein
distance~\cite{villani2009optimal}.

\begin{proposition}\label{prop:curv}
$\Delta\kappa\ge0$, and $\Delta\kappa=0$ whenever $E'_+\subseteq E$.
\end{proposition}
\begin{proof}
$W_1\ge0$. If $E'_+\subseteq E$, then $G\cup E'_+=G$, so the two empirical measures coincide and $W_1=0$.
\end{proof}
Thus a response that only restates GT edges incurs no distortion. Every novel claim edge raises endpoint degrees
and, unless it closes triangles, adds a negatively curved edge.

\paragraph{(3) Spectral divergence.}
Let $V'$ be the vertices touched by $E'_+$ and $H$ the \emph{$k$-hop shadow} of $G$ on $V'$: $(u,v)\in E(H)$ iff
$d_G(u,v)\le k$ (we use $k=2$, the retrieval depth). With
$L_{\mathrm{sym}}=I-D^{-1/2}AD^{-1/2}$,
\begin{equation}
D_{\mathrm{spec}}=\tfrac12\,W_1\big(\spec L_{\mathrm{sym}}(G'_+),\ \spec L_{\mathrm{sym}}(H)\big).
\end{equation}
\begin{lemma}\label{lem:spec}
$D_{\mathrm{spec}}\in[0,1]$, and $D_{\mathrm{spec}}=0$ if $G'_+$ and $H$ are isomorphic on $V'$.
\end{lemma}
\begin{proof}
All eigenvalues of $L_{\mathrm{sym}}$ lie in $[0,2]$~\cite{chung1997spectral}. Both spectra are uniform measures on
$|V'|$ points of $[0,2]$, so $W_1\le2$. Isomorphic graphs are cospectral.
\end{proof}
Whereas $\Delta\kappa$ is local (edge neighbourhoods), $D_{\mathrm{spec}}$ compares the global shape of the claim
graph with what the KG says the same entities look like.

\paragraph{(4) Bayesian provenance support.}
For each claim edge, let $s_e=g(u,v)$ if a GT edge exists, $s_e=\max_w g(u,w)g(w,v)$ if only two-hop paths exist (the
probability that both independent hops are genuine), and $s_e=0$ otherwise or if an endpoint is unaligned. With a
uniform prior on the support rate,
\begin{equation}
\pi_{\mathrm{prov}}=\E\big[\mathrm{Beta}\big(1+\textstyle\sum_e s_e,\ 1+\sum_e(1-s_e)\big)\big]=\frac{1+\sum_e s_e}{2+|E'_+|}.
\end{equation}

\paragraph{(5) Evidence similarity.}
$\sigma_{\mathrm{sim}}=\frac{1}{|C|}\sum_{c\in C}\max_{s\in\mathcal S}\cos(\phi(c),\phi(s))$ over verbalised claims $C$
and all corpus sentences $\mathcal S$. It compensates for KG recall: a true claim that Phase~1 missed can still find
its sentence.

\paragraph{(6, 7) NLI consistency.}
Embedding similarity measures topical closeness, not support: ``X uses Y'' and ``X does not use Y'' are near
neighbours. For each claim $c$ we retrieve its $k=3$ nearest corpus sentences $\mathcal S_c$ and score each pair with a
cross-encoder NLI model~\cite{he2023debertav3} fine-tuned on SNLI and MultiNLI~\cite{bowman2015snli,williams2018mnli}.
Following the zero-shot SummaC score~\cite{laban2022summac}, the consistency of a claim is
\begin{equation}
\gamma(c)=\max_{s\in\mathcal S_c}\big[P_{\mathrm{ent}}(s\Rightarrow c)-P_{\mathrm{con}}(s\Rightarrow c)\big]\in[-1,1].
\end{equation}
Subtracting \emph{per premise} before maximising matters: a separate maximum of contradiction lets any unrelated
sentence that happens to contradict a short claim dominate. We use the mean $\bar\gamma_{\mathrm{NLI}}$ over asserted
claims (how much of the answer is stated) and the minimum $\gamma^{\min}_{\mathrm{NLI}}$ (its weakest claim).

\paragraph{(8) Out-of-corpus novelty.}
Over the content stems $T$ of the asserted claims, with smoothed inverse document frequency
$\mathrm{idf}(t)=\log\frac{1+N}{1+\mathrm{df}(t)}+1$ over $N$ chunks and corpus vocabulary $\mathcal V$:
\begin{equation}
\nu_{\mathrm{ooc}}=\frac{\sum_{t\in T}\mathrm{idf}(t)\,\mathbb 1[t\notin\mathcal V]}{\sum_{t\in T}\mathrm{idf}(t)}.
\end{equation}
Restricting to asserted claims is essential. An answer that \emph{refutes} ``genetic algorithms'' mentions an
out-of-corpus term without asserting it.

\paragraph{(9) Contradiction rate.}
$\kappa_{\mathrm{con}}$ is the fraction of claims that (i)~are \textsc{refuted} by the answer although a grounded GT
edge exists, (ii)~give a second object to a subject under a relation marked functional in the KG, or (iii)~invert an
asymmetric relation (e.g.\ \textsc{outperforms}, \textsc{part\_of}, \textsc{is\_a}) whose reverse is in the KG.

\paragraph{(10) Entity alignment.}
$\rho_{\mathrm{align}}$ is the fraction of the answer's entities that align to a KG entity: an answer built on
entities the papers never mention is unlikely to be supported by them.

\paragraph{(11--13) Model-side signals.}
The verbalised confidence $c_{\mathrm{LLM}}$; the mean negative log-likelihood of the generated tokens,
$\bar u_{\mathrm{NLL}}=-\frac1L\sum_{\ell}\log p(w_\ell\mid w_{<\ell})$, an intrinsic uncertainty
signal~\cite{kadavath2022know,farquhar2024semantic} obtained from the same API call at no extra cost (zero when the
model does not expose token probabilities); and the hedging ratio $h_{\mathrm{spec}}$, the share of claims the answer
itself marks as speculative. None of them is trusted a priori: their direction and weight are learned.

\paragraph{Claim-level vectors.}
Besides the answer vector $\bm x\in\R^{13}$, every claim $c_i$ of an answer receives a vector $\bm z_i\in\R^{15}$:
provenance strength and support tier (direct edge, two-hop path), conflict flag, polarity, endpoint alignment,
NLI entailment, contradiction and consistency, embedding similarity to the nearest sentence, out-of-corpus rate, the
claim's own normalised sheaf residual $\log(1+\|\delta_e\|^2/\bar\delta^2)$ and the signed log Forman curvature of the
grafted edge. An answer is thus a \emph{bag} $\{\bm z_i\}$, which \cref{sec:phase4} exploits.

\paragraph{Hallucination extent.}
A binary verdict does not say \emph{how much} of an answer is unsupported. We therefore also report the claim-level
extent
\begin{equation}
\eta(a)=\frac{\big|\{e\in E'_+:\ \text{no grounded GT path, or } e \text{ is in conflict}\}\big|}{|E'_+|}\in[0,1],
\end{equation}
which is available for every answer, including before any calibration data exist. $\eta$ is a descriptive measure
reported alongside the calibrated probability $1-S_{\mathrm{final}}$; it is not an input to the learned scorer.

\paragraph{Local KG fidelity.}
The retrieved subgraph inherits its own reliability,
$\SKGq=\big(\sum_{e\in E}g_e+S_{\mathrm{KG}}\big)/(|E|+1)$, which is the mean grounding shrunk towards the global
score with one prior pseudo-edge.

\subsection{Phase 4: Learned Scoring and Conformal Decisions}
\label{sec:phase4}

\paragraph{Data and splits.}
Each labelled response contributes $(\bm x\in\R^{13},\{\bm z_i\}\subset\R^{15},\SKGq,y)$. A deterministic hash of the query identifier assigns 60\% /
20\% / 20\% to train, calibration and test. Because a query never changes split, models trained at different times
are compared on an identical test set without leakage.

\paragraph{Tournament.}
Ten families span the main inductive biases for small tabular data and for sets: logistic regression; gradient
boosting with \emph{monotone constraints} encoding the expected direction of each feature (\cref{tab:features});
random forest; extra trees; gradient boosting; RBF-SVM; MLP; a stacking ensemble; a feature-attention residual network;
and the multiple-instance network below. Hyper-parameters are tuned by inner 3-fold CV, and generalisation is
estimated by outer 5-fold CV to avoid selection bias~\cite{cawley2010overfitting}. The champion minimises
out-of-fold log-loss, a strictly proper scoring rule that rewards calibrated probabilities rather than ranking
alone~\cite{gneiting2007scoring}. It is refitted on the training split, and Platt scaling~\cite{platt1999probabilistic}
is fitted on its \emph{nested out-of-fold} logits, which are honest by construction, giving
$S_{\mathrm{query}}=\sigma(a\,\ell(\bm x)+b)$. For interpretability we also report the standardised logistic surrogate
$\logit P(y=1\mid\bm x)=b+\sum_j w_j z_j$, whose coefficients $w_j$ are the \emph{learned multipliers} of the features.

\paragraph{Claim-level localisation by multiple-instance learning.}
Labels exist only for whole answers, yet an answer is hallucinated if and only if \emph{at least one} of its claims is
unsupported: precisely the standard multiple-instance assumption~\cite{dietterich1997mil}. We therefore model an
answer as the bag $\{\bm z_i\}_{i=1}^{n}$ of its claim vectors. Each claim is encoded as $\bm h_i=f_\theta(\bm z_i)$ and
pooled with gated attention~\cite{ilse2018attention},
\begin{equation}
a_i=\operatorname{softmax}_i\!\Big(\bm w^\top\big[\tanh(V\bm h_i)\odot\operatorname{sigm}(U\bm h_i)\big]\Big),\qquad
\bm b=\sum_i a_i\bm h_i,
\end{equation}
and the answer logit is $\ell=g\big([\bm b;\,A\bm x]\big)$, so the network sees both the bag and the 13 answer-level
features. A second head gives each claim a probability $r_i=\operatorname{sigm}(\bm c^\top\bm h_i)$ of being
unsupported, tied to the answer label through a noisy-OR,
\begin{equation}
\mathcal L=\mathrm{BCE}\big(\sigma(\ell),y\big)+\lambda\,\mathrm{BCE}\Big(\textstyle\prod_i(1-r_i),\,y\Big),
\end{equation}
which encodes the multiple-instance assumption ($P(\text{supported})=\prod_i P(\text{claim } i \text{ supported})$ under
independence) and is what turns $r_i$ into a \emph{claim-level hallucination detector} trained without any
claim-level label. $\lambda$ and the width are tuned in the inner CV. If another family wins the tournament, the
network is still fitted as a companion localiser, so claim-level scores are always reported.

\paragraph{Learned log-linear pool.}
Earlier versions combined scores as $0.2\,S_{\mathrm{KG}}+0.8\,S_{\mathrm{query}}$. We instead learn
\begin{equation}
\logit S_{\mathrm{final}}=b_0+b_1\logit S_{\mathrm{query}}+b_2\logit\SKGq
\label{eq:pool}
\end{equation}
by maximum likelihood (Newton--Raphson with a weak $L_2$ prior) on \emph{out-of-fold} $S_{\mathrm{query}}$, so the pool
never sees in-sample over-confidence.

\begin{proposition}\label{prop:pool}
(a) The normalised geometric pool $p\propto S_{\mathrm{query}}^{\eta}S_{\mathrm{KG}}^{1-\eta}$ is the special case
$b_0=0$, $b_1=\eta$, $b_2=1-\eta$ of \cref{eq:pool}. (b) If the KG term is constant across queries, $b_2$ is not
identifiable.
\end{proposition}
\begin{proof}
(a) Normalising $p\propto\prod_i p_i^{w_i}$ over $\{0,1\}$ gives
$\frac{p}{1-p}=\prod_i\big(\frac{p_i}{1-p_i}\big)^{w_i}$, i.e.\ $\logit p=\sum_i w_i\logit p_i$. (b) If
$\logit S_{\mathrm{KG}}\equiv\lambda$, then $b_0+b_2\lambda$ enters the likelihood only through their sum, so any
$(b_0-\delta\lambda,\,b_2+\delta)$ yields the same likelihood.
\end{proof}
\Cref{prop:pool}(b) explains why a \emph{global} $S_{\mathrm{KG}}$ weight can never be learned from one corpus, and why
we pool the per-query $\SKGq$. If $\SKGq$ is nonetheless constant, the implementation drops $b_2$.

\paragraph{Mondrian conformal decision.}
On the calibration split ($n_c$ examples of class $c$), compute the non-conformity scores
$r_i=1-\hat p_{y_i}(x_i)$, where $\hat p_1=S_{\mathrm{final}}$ and $\hat p_0=1-S_{\mathrm{final}}$, separately per class.
Let $q_c$ be the $\lceil(n_c+1)(1-\alpha)\rceil$-th smallest score of class $c$ (or $1$ if that index exceeds $n_c$),
and set
\begin{equation}
\Gamma(x)=\{c\in\{0,1\}:\ 1-\hat p_c(x)\le q_c\}.
\end{equation}
\begin{proposition}[class-conditional validity~\cite{vovk2012conditional}]
If the calibration examples of class $c$ and a test example of class $c$ are exchangeable, then
$\Pr\big(c\in\Gamma(x)\mid y=c\big)\ge1-\alpha$.
\end{proposition}
\begin{proof}[Proof sketch]
Conditionally on $y=c$, the test score is exchangeable with the $n_c$ calibration scores. Its rank among the $n_c+1$
scores is therefore uniform, so it is at most the $\lceil(n_c+1)(1-\alpha)\rceil$-th smallest with probability at
least $1-\alpha$~\cite{lei2018distribution,angelopoulos2023gentle}.
\end{proof}
The \textsc{abstain} rate is the price of the guarantee. It shrinks as the scorer separates the classes, and its
magnitude is itself an evaluation metric. The only free parameter is the risk level $\alpha$.

\paragraph{Champion--challenger promotion.}
Retraining produces a challenger that replaces the production model only if its test log-loss is lower on the
\emph{same} test split. Unlabelled production traffic is logged for drift monitoring and is never self-labelled.

% =============================================================================================
\section{Experimental Setup}
\label{sec:setup}

\paragraph{Corpus and KG.}
The KG is built from DeepWalk, LINE and node2vec (\NDocs{} documents, \NChunks{} chunks). Of \NRawTriples{} extracted
triples, \NTriples{} pass evidence verification (\AcceptRate\%). They resolve to \NEntities{} entities, \NShared{} of
which are shared across papers. The fidelity pillars are $\Phi_{\mathrm{prov}}=\PhiProv\%$,
$\Phi_{\mathrm{cov}}=\PhiCov\%$, $\Phi_{\mathrm{topo}}=\PhiTopo\%$ and $\Phi_{\mathrm{cons}}=\PhiCons\%$, giving
$S_{\mathrm{KG}}=\SKG\%$. On a uniform sample of $n=\KGAuditN$ triples checked by \KGAuditor, strict precision is
\KGPrec\% (95\% Wilson CI [\KGPrecLo\%, \KGPrecHi\%]) and lenient precision \KGPrecLenient\% (\cref{tab:kgprec}).

\paragraph{Benchmark.}
We generate \NItems{} questions from the corpus in four categories: \textsc{factual} (30\%, answerable from one
chunk), \textsc{cross-doc} (10\%, requiring chunks from two papers), \textsc{borderline} (20\%, on-topic but not
answerable from the papers) and \textsc{false-premise} (40\%, presupposing something a chunk contradicts). Each item
stores a gold evidence quote that must fuzzy-match its source chunk, and items without one are discarded. Premises are
phrased naturally, and questions beginning with ``Assume'' are rejected so a detector cannot learn a template.
Near-duplicates are removed (cosine $\ge0.92$), and papers are balanced within categories.

\paragraph{Response-level labels.}
An LLM judge~\cite{zheng2023judging} receives the question and its type, the gold reference (answer, premise
correction or reason for unanswerability), the gold chunk(s), six retrieved corpus passages and the answer. It
returns \textsc{supported}, \textsc{partially supported} or \textsc{unsupported}, and $y=1$ iff the verdict is
\textsc{supported}. The primary judge is \texttt{\PrimaryJudge}. Because it sees the gold evidence it verifies
claims against a source rather than grading its own fluency; to quantify residual judge error, a second, stronger judge
(\texttt{\CrossJudge}) independently re-labels a fixed, category-stratified sample of \CrossJudgeN{} responses with the
identical prompt, giving inter-judge agreement $\kappa=\CrossJudgeKappa$ (raw agreement \CrossJudgeAgree\%).
Training uses the primary labels only. On a stratified sample of \JudgeN{} responses annotated by the authors,
binary human--judge agreement is Cohen's $\kappa=\JudgeKappa$.
\todo{report judge model and annotator protocol}. Overall, \SupportedRate\% of responses are supported.

\paragraph{Protocol and metrics.}
The splits contain \NTrain{} / \NCal{} / \NTest{} responses. We report AUROC, log-loss, Brier score and expected
calibration error (ECE, 10 bins) of $S_{\mathrm{final}}$ on the test split, together with empirical coverage (overall
and per class), abstention rate and selective accuracy (accuracy on non-abstained decisions) at $\alpha=\Alpha$.
\textbf{Uncertainty and tests:} every test metric carries a 95\% percentile bootstrap interval
(1{,}000 resamples)~\cite{efron1993bootstrap}, and AUROC differences are tested with the paired DeLong
test~\cite{delong1988comparing,sun2014fast}. \textbf{Baselines:} the LLM's verbalised confidence $c_{\mathrm{LLM}}$;
every single feature with its expected sign; and a \emph{text-only NLI checker}, the same learner on
$\{\sigma_{\mathrm{sim}},\bar\gamma_{\mathrm{NLI}},\gamma^{\min}_{\mathrm{NLI}},\nu_{\mathrm{ooc}}\}$, which stands for the
literature's NLI-based verifiers~\cite{laban2022summac,min2023factscore}. \textbf{Ablations:} the same standardised
logistic learner on each family (topological, evidential, model-side), on the graph-only features, and on all
features without the topological or without the model-side family.

\paragraph{Implementation.}
Embeddings come from \texttt{all-MiniLM-L6-v2}~\cite{reimers2019sbert}. Graph computations use NetworkX and SciPy,
learning uses scikit-learn~\cite{pedregosa2011sklearn} and PyTorch, and the graph is stored in Neo4j AuraDB with a
local JSON mirror. API calls run in an I/O thread pool with disk caching (optionally through a batch API).
CPU-bound stages run single-process with two BLAS threads at reduced priority. The full pipeline is resumable from
checkpoints. The system is exposed as a REST service with an interactive web interface for corpus ingestion, single
and batch auditing, benchmark execution and inspection of the calibration layer; all results are persisted per
workspace and restored on reload.

% =============================================================================================
\section{Results}
\label{sec:results}

\TableLeaderboard
\TableAblation

\paragraph{Detection and calibration.}
On the held-out test split ($n=\NTest$), the champion (\Champion) yields $S_{\mathrm{final}}$ with AUROC \TestAUC,
log-loss \TestLL, Brier \TestBrier{} and ECE \TestECE. The LLM's verbalised confidence reaches AUROC \LLMAUC{} with
ECE \LLMECE{} (\cref{fig:roc,fig:rel}). The 95\% bootstrap interval of the audit's AUROC is
[\TestAUCLo, \TestAUCHi] against [\LLMAUCLo, \LLMAUCHi] for the LLM's confidence; the paired DeLong test gives
$\Delta\mathrm{AUC}=\DeltaLLM$ ($p$ = \PvalLLM). Against the text-only NLI checker the difference is
$\Delta\mathrm{AUC}=\DeltaText$ ($p$ = \PvalText). \Cref{tab:leaderboard} lists the out-of-fold tournament.
% TODO(after run): interpret the gap between S_final and c_LLM; note which families are competitive
% and whether the monotone-constrained model is within noise of the champion.

\TableKgprec
\TableTripleweights
\TableSens

\paragraph{How much is unsupported.}
The claim-level extent $\eta$ averages \ExtentHal\% on answers the judge labels hallucinated and \ExtentSup\% on
supported answers; used alone as a score ($1-\eta$), it reaches AUROC \ExtentAUC.
% TODO(after run): comment on cases where eta is high but the answer is supported (KG recall gaps).

\paragraph{What each family contributes.}
\Cref{tab:ablation} isolates the feature families, and \cref{tab:features} reports single-feature AUROC, the learned
weights $w_j$ and permutation importance (\cref{fig:imp}). Whether the topological family adds information
beyond the evidential and model-side features is tested directly: all features versus all-but-topological gives
$\Delta\mathrm{AUC}=\DeltaTopo$ ($p$ = \PvalTopo).
% TODO(after run): discuss (i) topological-only vs evidential-only, (ii) the drop when removing the
% topological family (the paper's central hypothesis), (iii) the weight learned for c_LLM.

\TableFeatures

\paragraph{Learned pool.}
The fitted pool is $\logit S_{\mathrm{final}}=\PoolBzero+\PoolBq\logit S_{\mathrm{query}}+\PoolBkg\logit\SKGq$.
% TODO(after run): compare b_query / b_kg with the discarded fixed 0.8 / 0.2 split.

\paragraph{Conformal decisions.}
At $\alpha=\Alpha$, empirical coverage on the test split is \Coverage\% overall, \CoverageSup\% on supported and
\CoverageHal\% on hallucinated responses, with an abstention rate of \AbstainRate\% and selective accuracy of
\SelAcc\%. For display, the implied score cut-offs are $\tau_{\mathrm{ver}}=\TauVer$ and $\tau_{\mathrm{hal}}=\TauHal$;
these are \emph{outputs} of calibration, not inputs. \Cref{tab:percat} breaks the decisions down by category.
% TODO(after run): comment on false-premise answers that correctly reject the premise vs. those that accept it.

\TablePercat

\begin{figure}[t]\centering
\IfFileExists{figures/roc.pdf}{\includegraphics[width=0.92\columnwidth]{figures/roc.pdf}}{\fbox{\parbox{0.85\columnwidth}{\centering ROC figure pending: run \texttt{scripts/export\_paper\_assets.py}}}}
\caption{ROC on held-out responses: learned audit score $S_{\mathrm{final}}$, pre-pool $S_{\mathrm{query}}$ and the
LLM's verbalised confidence.}\label{fig:roc}
\end{figure}

\begin{figure}[t]\centering
\IfFileExists{figures/reliability.pdf}{\includegraphics[width=0.92\columnwidth]{figures/reliability.pdf}}{\fbox{\parbox{0.85\columnwidth}{\centering Reliability figure pending}}}
\caption{Reliability diagram (10 equal-width bins). The diagonal is perfect calibration.}\label{fig:rel}
\end{figure}

\begin{figure}[t]\centering
\IfFileExists{figures/importance.pdf}{\includegraphics[width=0.92\columnwidth]{figures/importance.pdf}}{\fbox{\parbox{0.85\columnwidth}{\centering Importance figure pending}}}
\caption{Permutation importance of the champion: increase in test log-loss when a feature is shuffled (mean $\pm$
s.d., 20 repeats).}\label{fig:imp}
\end{figure}

% =============================================================================================
\section{Discussion and Limitations}
\label{sec:limits}

\paragraph{What the guarantee does and does not say.}
Coverage holds marginally within each class and assumes exchangeability between calibration and deployment
responses. A change of corpus, audited model, prompt or question distribution breaks this assumption. The
champion--challenger loop with fresh labels is the remedy, and unlabelled telemetry is retained to detect drift.

\paragraph{Labels.}
Response-level labels from an LLM judge are scalable but imperfect. We therefore report agreement with human
annotation and with a second, stronger judge, and judge bias remains a threat to validity. When the primary judge
shares a model family with the audited model, self-preference is possible; giving the judge the gold evidence and
reporting inter-judge agreement (\S\ref{sec:setup}) bound its effect.

\paragraph{KG recall.}
A true claim absent from the KG looks unsupported to the graph features. The text-level features
$\sigma_{\mathrm{ent}}$ and $\nu_{\mathrm{ooc}}$, and the learned weighting, partially compensate. $S_{\mathrm{KG}}$ and
$\SKGq$ make the remaining dependence explicit rather than hidden.

\paragraph{Scope.}
Three papers form a narrow, homogeneous domain. Entity-name embeddings are a coarse stalk for the sheaf. Richer
restriction maps (e.g.\ learned orthogonal maps per relation~\cite{gebhart2023knowledge}) and larger, multi-domain
corpora are natural next steps.

\paragraph{Ethics.}
The corpus consists of published papers used for research reproducibility. The system is designed to \emph{flag}
unsupported content for human review, and the \textsc{abstain} outcome exists to avoid over-confident automated
judgements.

% =============================================================================================
\section{Conclusion}
\label{sec:conclusion}

EpistemeGraph treats hallucination auditing as calibrated measurement. It builds a fallible but quantified reference
graph, extracts the claim graph of an answer independently of its generation, measures the mismatch with sheaf,
curvature, spectral and evidential instruments, and learns, rather than prescribes, how these measurements become a
decision with a stated error rate. The full system, benchmark, prompts and every intermediate artefact are released
to make each number in this paper reproducible.

\bibliographystyle{plainnat}
\bibliography{references}

% =============================================================================================
\appendix
\onecolumn
\section{Algorithms}
\label{app:alg}

\begin{algorithm}[h]
\caption{Audit of one query (inference)}
\begin{algorithmic}[1]
\Require query $q$; KG snapshot; production scorer $\hat p$; pool $(b_0,b_1,b_2)$; quantiles $(q_0,q_1)$
\State $a, c_{\mathrm{LLM}} \gets \textsc{Answer}(q)$ \Comment{stage A}
\State $G'\gets\textsc{ExtractClaims}(q,a;\ \text{KG relation vocabulary, nearest KG entities})$ \Comment{stage B}
\State $G\gets\textsc{Retrieve}(\text{entities of }G')$; align $V'\to V$
\State $x\gets(\mathcal E_{\mathcal F},\Delta\kappa,D_{\mathrm{spec}},\pi_{\mathrm{prov}},\sigma_{\mathrm{ent}},\nu_{\mathrm{ooc}},\kappa_{\mathrm{con}},c_{\mathrm{LLM}})$
\State $S_{\mathrm{query}}\gets\hat p(x)$;\quad $S_{\mathrm{final}}\gets\sigma\big(b_0+b_1\logit S_{\mathrm{query}}+b_2\logit\SKGq\big)$
\State $\Gamma\gets\{c:1-\hat p_c\le q_c\}$ with $\hat p_1=S_{\mathrm{final}}$, $\hat p_0=1-S_{\mathrm{final}}$
\State \Return \textsc{verified} if $\Gamma=\{1\}$; \textsc{hallucinated} if $\Gamma=\{0\}$; else \textsc{abstain}
\end{algorithmic}
\end{algorithm}

\begin{algorithm}[h]
\caption{Phase 4 training}
\begin{algorithmic}[1]
\Require labelled set $\{(x_i,\SKGq_i,y_i)\}$, risk level $\alpha$
\State split by $\mathrm{hash}(\text{query id})$ into train / cal / test
\For{each model family $m$ in the tournament}
  \State OOF$_m\gets$ outer 5-fold CV of (inner 3-fold grid search of $m$) on train
\EndFor
\State $m^\star\gets\argmin_m \mathrm{LogLoss}(\text{OOF}_m)$; fit Platt-calibrated $m^\star$ on train
\State fit pool \cref{eq:pool} on out-of-fold calibrated $S_{\mathrm{query}}$ of train
\State $q_c\gets$ class-$c$ conformal quantile of $1-\hat p_{y}(x)$ on cal
\State evaluate on test; promote if test log-loss improves on production
\end{algorithmic}
\end{algorithm}

\section{Incremental Neo4j schema}
\label{app:cypher}
\begin{footnotesize}
\begin{verbatim}
CREATE CONSTRAINT entity_key_unique IF NOT EXISTS FOR (e:Entity) REQUIRE e.name_key IS UNIQUE;

UNWIND $batch AS row
MERGE (e:Entity {name_key: row.name_key})
  ON CREATE SET e.name = row.name, e.type = row.type, e.aliases = row.aliases, e.source_pdfs = row.source_pdfs
  ON MATCH  SET e.aliases     = coalesce(e.aliases,[]) + [a IN row.aliases WHERE NOT a IN coalesce(e.aliases,[])],
                e.source_pdfs = coalesce(e.source_pdfs,[]) + [p IN row.source_pdfs WHERE NOT p IN coalesce(e.source_pdfs,[])];

UNWIND $batch AS row
MATCH (u:Entity {name_key: row.subj_key}), (v:Entity {name_key: row.obj_key})
MERGE (u)-[r:RELATED_TO {type: row.rel}]->(v)
WITH r, row, (r.grounding_score IS NULL OR row.grounding >= r.grounding_score) AS better
SET r.grounding_score = CASE WHEN better THEN row.grounding  ELSE r.grounding_score END,
    r.evidence_quote  = CASE WHEN better THEN row.best_quote ELSE r.evidence_quote  END,
    r.chunk_ids       = coalesce(r.chunk_ids,[]) + [c IN row.chunk_ids WHERE NOT c IN coalesce(r.chunk_ids,[])]
SET r.support_count = size(r.chunk_ids);
\end{verbatim}
\end{footnotesize}

\section{Prompts}
\label{app:prompts}
Full prompts are versioned in the repository: KG extraction (\texttt{kg\_builder.py}, \texttt{PROMPT\_VERSION}),
answering and claim extraction (\texttt{llm\_interceptor.py}), benchmark generation
(\texttt{scripts/generate\_benchmark.py}) and judging (\texttt{scripts/run\_benchmark.py}).

\end{document}
